%% kapa: exam=1 round=34 year=2023 period=1 form=A subject=civil_law subject_ko=민법 seq=34 number=34
%% source: Q-Net 감정평가사 기출문제 (cst00309) artlSeq=5211789 file=raw/1008/5211789/2023년도 제34회 감정평가사 1차시험 1교시.pdf
%% conversion: kordoc 4.24.51 -> tools/gen1.py (auto); review=proofread-sonnet; flags=none
\begin{question}{34}
지역권에 관한 설명으로 옳지 않은 것은? (다툼이 있으면 판례에 따름)
\begin{choices}
  \item 지역권은 요역지의 사용가치를 높이기 위해 승역지를 이용하는 것을 내용으로 하는 물권이다.
  \item 요역지와 승역지는 서로 인접한 토지가 아니어도 된다.
  \item 요역지 공유자 중 1인에 대한 지역권 소멸시효의 정지는 다른 공유자를 위하여도 효력이 있다.
  \item 지역권자는 승역지의 점유침탈이 있는 경우, 지역권에 기하여 승역지 반환청구권을 행사할 수 있다.
  \item 지역권은 계속되고 표현된 것에 한하여 시효취득할 수 있다.
\end{choices}
\end{question}
\answer{4}
